% USERS_MANUAL_JA.tex — Xinu (RAZ 拡張版) ユーザズマニュアル (日本語版)
% Build: lualatex USERS_MANUAL_JA.tex   (目次のため 2 回実行)
\documentclass[a4paper,11pt]{ltjsarticle}

\usepackage{luatexja-fontspec}
\setmonofont{Menlo}[Scale=0.85]
\usepackage[a4paper,margin=2.0cm]{geometry}
\usepackage{amssymb}
\usepackage{xcolor}
\usepackage{fancyvrb}
\usepackage{longtable}
\usepackage{booktabs}
\usepackage{array}
\usepackage{enumitem}
\usepackage{newunicodechar}
\usepackage[hidelinks]{hyperref}

\newunicodechar{→}{\textrightarrow}
\newunicodechar{←}{\textleftarrow}
\newunicodechar{↑}{\textuparrow}
\newunicodechar{↓}{\textdownarrow}
\newunicodechar{↔}{\ensuremath{\leftrightarrow}}
\newunicodechar{×}{\ensuremath{\times}}
\newunicodechar{…}{\dots}
\newunicodechar{≤}{\ensuremath{\leq}}
\newunicodechar{≥}{\ensuremath{\geq}}
\newunicodechar{≈}{\ensuremath{\approx}}
\newunicodechar{≒}{\ensuremath{\fallingdotseq}}
\newunicodechar{°}{\textdegree}

\DefineVerbatimEnvironment{code}{Verbatim}%
  {fontsize=\small,frame=single,framerule=0.3pt,rulecolor=\color{gray!60},%
   framesep=5pt,xleftmargin=2pt,xrightmargin=2pt}

\setlength{\parindent}{0pt}
\setlength{\parskip}{4pt}
\renewcommand{\arraystretch}{1.25}

\title{\textbf{Xinu (RAZ 拡張版) ユーザズマニュアル}\\[3pt]
  \large XFS ファイルシステム・\texttt{cc} C コンパイラ・%
  \texttt{abclc} アクタ言語・\texttt{wm} ウィンドウマネージャ}
\author{}
\date{}

\begin{document}
\maketitle
\vspace{-2.0em}
\begin{center}\rule{0.9\textwidth}{0.4pt}\end{center}

本マニュアルは、Embedded Xinu をベースに、階層型ファイルシステム \textbf{XFS}
(RAM ディスク上)、組込みのちっちゃな C コンパイラ \textbf{cc} とバイトコード VM
\textbf{a.out}、アクタ指向小型言語 \textbf{ABCL/c+} とその C トランスレータ
\textbf{abclc}、ウィンドウマネージャ \textbf{wm} (ARM/QEMU 上のフレームバッファ
GUI)、Unix 風シェル拡張 (\texttt{ls}/\texttt{cd}/\texttt{cat}/\texttt{make}/%
\texttt{edit} など) を組み込んだ拡張版 Xinu の使い方をまとめたものです。本家
Embedded Xinu の機能 (スレッド・セマフォ・ネットワーク等) はそのまま使えます。

{\small\tableofcontents}
\bigskip

%======================================================================
\section{はじめに}
Xinu は Douglas Comer が教育用に開発した小型 UNIX 風 OS です。本拡張版はこれに
「OS 上で C/ABCL ソースを編集・コンパイル・実行できる自己完結したワークベンチ」を
載せたもので、ホスト側のクロスコンパイラに頼らずとも、QEMU の中だけでプログラム
開発のループを完結できます。

\begin{center}\footnotesize
\begin{tabular}{@{}l l p{6.6cm}@{}}
\toprule
\textbf{プラットフォーム名} & \textbf{用途} & \textbf{特徴}\\
\midrule
\verb|arm-qemu| & QEMU \verb|-M versatilepb| & UART シェル + フレームバッファ + WM 自動起動\\
\verb|arm-rpi|  & Raspberry Pi 実機 & UART / HDMI / USB キーボード\\
\bottomrule
\end{tabular}
\end{center}
XFS / cc / abclc / make / edit はどのプラットフォームでも同じように使えます。
WM 関連 (\verb|rotlines|, \verb|wm_line| など) は arm-qemu でのみ有効です。

%======================================================================
\section{起動方法}
\subsection{ビルド}
\begin{code}
cd compile
make PLATFORM=arm-qemu        # 既定。arm-rpi も可
\end{code}
成功すると \verb|compile/xinu.boot| が生成されます。

\subsection{QEMU で実行}
ホストが macOS/Linux で \verb|qemu-system-arm| がインストール済みであれば、用意
してあるラッパスクリプトが便利です。

\textbf{コンソール専用モード (フレームバッファなし):}
\begin{code}
./compile/run-console.sh
\end{code}
UART シェル (\verb|xsh$|) が端末に多重化されます。終了は \textbf{Ctrl-A
\textrightarrow{} x}。環境変数: \verb|XINU_MEM| (既定 \verb|128M|)、
\verb|XINU_PLATFORM| (既定 \verb|arm-qemu|)。

\textbf{ウィンドウモード (WM デモ込み):}
\begin{code}
./compile/run-window.sh
\end{code}
LCD フレームバッファが macOS の Cocoa ウィンドウに表示されます。UART シェルは
起動した端末で動きます (\verb|-serial mon:stdio| で QEMU モニタも同居)。終了は
ウィンドウを閉じるか端末で \textbf{Ctrl-A \textrightarrow{} x}。ウィンドウは角を
ドラッグで最大 4 倍まで滑らかに拡大可能 (zoom-to-fit)。

\subsection{起動シーケンス}
\verb|system/main.c| の \verb|main()| スレッドが以下を順に行います。
\begin{enumerate}[leftmargin=1.6em,itemsep=1pt]
\item OS バージョン・メモリレイアウトを表示
\item \verb|xfsBootstrap()| で RAMDISK0 を XFS でフォーマット \& \verb|/| にマウント
\item \verb|/home/hello.c|, \verb|sum.c|, \verb|PingPong.abcl|,
  \verb|RotLines.abcl| をシード書き込み
\item ネットワークデバイスを open (\verb|NETHER| 有効時)
\item CONSOLE/TTY1 を open
\item arm-qemu のときだけ \verb|wm_main| スレッドを \verb|create()|
\item シェルスレッドを CONSOLE/TTY1 ごとに起動
\end{enumerate}
ブート完了後、シェルプロンプトに \verb|xsh$| と表示されればコマンド入力可能です。

%======================================================================
\section{シェルの使い方}
\subsection{基本}
\begin{code}
コマンド名 [引数...]    [< 入力リダイレクト] [> 出力リダイレクト] [&]
\end{code}
\verb|&| でバックグラウンド起動、\verb|<| \verb|>| で入出力デバイスへリダイレクト。
補完・履歴はありません。

\subsection{コマンド一覧 (本拡張版で追加・拡張されたもの)}
\begin{center}\footnotesize
\begin{longtable}{@{}l l p{7.2cm}@{}}
\toprule
\textbf{コマンド} & \textbf{使い方} & \textbf{説明}\\
\midrule \endhead
\verb|pwd|   & \verb|pwd|              & カレントディレクトリ表示\\
\verb|cd|    & \verb|cd [PATH]|        & ディレクトリ変更 (省略時は \verb|/|)\\
\verb|ls|    & \verb|ls [-l] [PATH]|   & 一覧。\verb|-l| で type/size 表示\\
\verb|cat|   & \verb|cat FILE...|      & ファイル内容表示\\
\verb|mkdir| & \verb|mkdir DIR...|     & ディレクトリ作成\\
\verb|rmdir| & \verb|rmdir DIR...|     & 空ディレクトリ削除\\
\verb|touch| & \verb|touch FILE...|    & 空ファイル作成 / mtime 更新\\
\verb|rm|    & \verb|rm FILE...|       & ファイル削除\\
\verb|cp|    & \verb|cp SRC DST|       & コピー (DST は truncate)\\
\verb|mv|    & \verb|mv SRC DST|       & リネーム/移動\\
\verb|write| & \verb|write FILE TEXT...| & 引数を空白区切りで連結し改行付きで書き込み\\
\verb|edit|  & \verb|edit FILE|        & 組込み Emacs 風エディタ\\
\verb|mkfs|  & \verb|mkfs DEVICE [VOL]| & ブロックデバイスを XFS でフォーマット\\
\verb|mount| & \verb|mount DEV MNT|    & XFS をマウント\\
\verb|umount|& \verb|umount MNT|       & XFS をアンマウント\\
\verb|cc|    & \verb|cc SRC [-o OUT]|  & C ソースを a.out にコンパイル\\
\verb|abclc| & \verb|abclc SRC.abcl [-o OUT]| & ABCL \textrightarrow{} C \textrightarrow{} a.out 一気通貫\\
\verb|make|  & \verb|make [TARGET...]| & カレントの Makefile を実行\\
\verb|run|   & \verb|run PATH|         & a.out を実行\\
\verb|rotlines| & \verb|rotlines [F] [DEG/F]| & WM 上に回転線を描く\\
\verb|clear| & \verb|clear|            & 画面クリア\\
\verb|sleep| & \verb|sleep N|          & N ms スリープ\\
\bottomrule
\end{longtable}
\end{center}

\textbf{暗黙の \texttt{run}。} シェルが知らないコマンド名を打つと、最後に
\verb|aoutRun(<その名前>)| を試します (\verb|shell/shell.c:332| 付近)。つまり
\verb|cc hello.c -o hello| のあとは \verb|xsh$ hello| でファイルパス相当の名前から
直接プログラムを起動できます。

%======================================================================
\section{XFS --- ファイルシステム}
\verb|include/xfs.h| で定義される 4 KB ブロックの階層型 FS で、起動時に 16 MB の
RAM ディスク (\verb|RAMDISK0|) 上に作られ \verb|/| にマウントされます。
\begin{center}\footnotesize
\begin{tabular}{@{}l l@{}}
\toprule \textbf{項目} & \textbf{値}\\ \midrule
ブロックサイズ & 4096 B\\
最大ファイル名長 & 56 B\\
inode サイズ & 128 B\\
直接ブロック & 12 (= 48 KB)\\
一段間接 & 1024 ブロック (≒ 4 MB)\\
二段間接 & あり\\
マジック & \verb|XFS!| (0x58465321)\\
\bottomrule
\end{tabular}
\end{center}
ディレクトリは inode の中身が \verb|struct xdirent[]| (1 エントリ 64 B) で表現
され、\verb|xfsReaddir()| で順次読み出します。\verb|xfsChdir()|/\verb|xfsGetcwd()|
は\textbf{スレッドごと}の絶対パス CWD を保持します (シェル A で \verb|cd /home|
してもシェル B には影響しません)。\verb|xfsBootstrap()| は \verb|RAMDISK0| 先頭
64 KB に XFS マジックがあるか調べ、なければ \verb|xfsMkfs| で初期化し \verb|/| に
マウントします。
\begin{code}
int fd = xfsOpen("/home/foo", XFS_O_RDWR | XFS_O_CREAT | XFS_O_TRUNC);
xfsWrite(fd, buf, n);  xfsRead(fd, buf, n);
xfsSeek(fd, 0, XFS_SEEK_SET);  xfsClose(fd);
xfsMkdir("/home/sub");  xfsRename("/home/old", "/home/new");  xfsUnlink("/home/foo");
\end{code}
RAMDISK0 は名前のとおり RAM です。\textbf{QEMU を終了するとファイルは消えます。}
永続的に残したいソースはホストでバックアップしてください。

%======================================================================
\section{cc / run --- C コンパイラとランチャ}
\subsection{サポートする C のサブセット}
\verb|system/cc.c| 冒頭に明記: 1 ソースファイル、\verb|int main(void)| 1 個 + 引数
なし補助関数。宣言は \verb|int x;| / \verb|int x = expr;| のみ。文は
\verb|if/else|, \verb|while|, \verb|for|, \verb|return|, ブロック, 式文。式は
整数・文字列・識別子・\verb|+ - * / %|・比較・\verb|! && |||・単項 \verb|-|・括弧。
組込み関数のみ呼出可。\verb|#include| は受理して無視。\textbf{未サポート:}
構造体・ポインタ・配列・浮動小数・複数ファイル。

\subsection{組込み関数 (Built-in)}
\begin{center}\footnotesize
\begin{tabular}{@{}l l p{6.6cm}@{}}
\toprule \textbf{BI} & \textbf{関数} & \textbf{用途}\\ \midrule
0 & \verb|printf(fmt, ...)| & 文字列・整数フォーマット\\
1 & \verb|puts(s)| & 文字列+改行\\
2 & \verb|putchar(c)| & 1 文字出力\\
3 & \verb|getchar()| & 1 文字入力\\
4 & \verb|exit(code)| & 終了\\
5 & \verb|rgb(r,g,b)| & BGR565 16bit カラー値\\
6 & \verb|wm_line(idx,x1,y1,x2,y2,color)| & WM のユーザ線スロット (0..7) に登録\\
7 & \verb|wm_render(on)| & ユーザ線レンダの有効/無効\\
8 & \verb|wm_clear()| & ユーザ線クリア\\
9 & \verb|sleep_ms(ms)| & スリープ\\
10/11 & \verb|isin(deg)|/\verb|icos(deg)| & Q12 固定小数 (4096 = 1.0)\\
12/13 & \verb|screen_w()|/\verb|screen_h()| & 画面サイズ取得\\
\bottomrule
\end{tabular}
\end{center}

\subsection{a.out 形式とシェルでの使い方}
\verb|include/aout.h| の \verb|struct aout_header|:
\begin{code}
+--------------------------+
| magic = "XAOU"           |
| version = 1              |
| code_size / const_size   |
| entry / nlocals          |
+--------------------------+
| バイトコード本体           |
| 文字列定数プール (零終端)  |
+--------------------------+
\end{code}
VM はスタック型 (256 entries) で、\verb|OP_PUSH_I32| \verb|OP_LOAD_LOC|
\verb|OP_ADD| \verb|OP_JZ| \verb|OP_CALL_BI| などの 1 バイトオペコードを持ちます
(\verb|system/aout.c| の \verb|aoutRun()|)。
\begin{code}
xsh$ cc /home/hello.c -o /home/hello
xsh$ run /home/hello
hello...
xsh$ /home/hello              # 暗黙 run でも OK
\end{code}

%======================================================================
\section{abclc --- ABCL/c+ アクタ言語}
ABCL/c+ は「クラス + メッセージ送信」によるアクタ指向の小型言語で、本実装では
\textbf{abclc が C へ変換し、その C を cc がバイトコード化}する 2 段階構成です。
\begin{code}
class ClassName {
    var field1;
    method methodName(p1, p2) {
        field1 = p1 + 1;
        send other.methodName(arg1, arg2);
    }
}
main {
    new ClassName a;  new ClassName b;
    send a.methodName(b, 10);
}
\end{code}
暗黙の識別子 \verb|self| (現在のアクタ ID)・\verb|sender| (送信元 ID)。1 メソッド
あたり \verb|send| は最大 1 回 (末尾呼び出し型) なので、1 アクタは 1 メッセージずつ
逐次処理されます。値はすべて \verb|int| 相当。制限: 4 クラス、1 クラス 16 メソッド、
8 フィールド、8 引数。
\begin{code}
xsh$ cd /home/abclcp/abclc
xsh$ abclc PingPong.abcl
abclc: translated PingPong.abcl -> PingPong.c
abclc: compiled PingPong.c -> PingPong
xsh$ run PingPong
\end{code}
\verb|-o NAME| を付けると \verb|NAME.c| (中間 C) と \verb|NAME| (a.out) の両方が
作られます。abclc が出力する C は cc のビルトインと少数のスケジューリング用ヘルパ
だけを使い、1 メッセージ = 1 タイムスライスで実行され、\verb|send| 後はキューに
ぶら下がります。

%======================================================================
\section{make --- 簡易ビルドシステム}
\verb|shell/xsh_make.c|: CWD に \verb|Makefile| がなければ \verb|*.c| と
\verb|*.abcl| を走査して自動生成し、\verb|TARGETS = ...| で指定された (または
ルール宣言順の) 各ターゲットをビルド。\verb|.c|\,\textrightarrow{}\,%
\verb|cc SRC -o TARGET|、\verb|.abcl|\,\textrightarrow{}\,\verb|abclc SRC -o TARGET|。
\begin{code}
# コメント
TARGETS = hello sum PingPong RotLines
hello:    hello.c
sum:      sum.c
PingPong: PingPong.abcl
RotLines: RotLines.abcl
\end{code}
明示ルールが無いターゲットは \verb|TARGET.c| \textrightarrow{} \verb|TARGET.abcl|
の順に存在チェックして自動推論します。
\begin{code}
xsh$ cd /home
xsh$ make
make: generated Makefile (2 .c, 0 .abcl)
    cc hello.c -o hello
    cc sum.c -o sum
make: built 2 target(s)
xsh$ ./hello
hello...
\end{code}

%======================================================================
\section{edit --- 組込みエディタ}
\verb|shell/xsh_edit.c| の Emacs 風モードレスエディタ (最大 400 行 × 256 文字)。
\verb|edit /home/hello.c|。存在しないパスは新規バッファで開きます。
\begin{center}\footnotesize
\begin{tabular}{@{}l l@{}}
\toprule \textbf{キー} & \textbf{動作}\\ \midrule
\verb|C-f| / \verb|C-b| & 1 文字 進む / 戻る\\
\verb|C-n| / \verb|C-p| & 1 行 下 / 上\\
\verb|C-a| / \verb|C-e| & 行頭 / 行末\\
←↑→↓ & 矢印 (ESC[ シーケンス)\\
\verb|C-d| & カーソル位置を削除\\
\verb|Backspace| & 1 文字戻して削除\\
\verb|Enter| & 改行挿入\\
\verb|Tab| & 空白 2 個\\
\verb|C-k| & 行末まで kill\\
\verb|C-l| & 再描画\\
\verb|C-x C-s| & 保存\\
\verb|C-x C-c| & 保存して終了\\
\verb|C-g| & 保存せず終了\\
\bottomrule
\end{tabular}
\end{center}
ステータス行に \verb|L行:C列  status| が表示されます。

%======================================================================
\section{wm --- ウィンドウマネージャ}
\verb|apps/wm.c| は arm-qemu の VersatilePB 上で PL110 LCD コントローラ
(\verb|0x10120000|, 16bpp BGR565) と PL050 PS/2 マウス (\verb|0x10007000|) を直接
叩いて、1024×1024 仮想画面 (実画面 640×480) のデスクトップ風 GUI を表示します。
\verb|run-window.sh| で起動するとブート時に自動で \verb|wm_main| が立ち上がります。
提供するもの: ステータスバー、ドラッグ可能なデモウィンドウ 3 枚、マウスカーソル、
ユーザ用描画フック \verb|wm_set_user_render(...)|、直接描画 API
(\verb|put_pixel_pub| / \verb|fill_rect_pub| / \verb|rect_outline_pub| /
\verb|draw_string_pub| / \verb|wm_draw_line|)。

C プログラム / ABCL からは a.out のビルトイン経由で \verb|rgb(r,g,b)|・
\verb|wm_render(1)|・\verb|wm_line(idx,...)|・\verb|wm_clear()|・
\verb|screen_w()|/\verb|screen_h()| を使うのが基本です (RotLines.abcl 参照)。
シェルコマンド \verb|rotlines| はもう少し低レベルで、\verb|wm_set_user_render|
で毎フレームの描画関数を直接 WM に登録します。

%======================================================================
\section{サンプルプログラム詳説}
ブート時に \verb|system/main.c| が以下 4 本を XFS にシードします。
\begin{center}\footnotesize
\begin{tabular}{@{}l l l@{}}
\toprule \textbf{パス} & \textbf{種類} & \textbf{主役}\\ \midrule
\verb|/home/hello.c| & C & \verb|printf|\\
\verb|/home/sum.c|   & C & \verb|for| ループ\\
\verb|.../PingPong.abcl| & ABCL & アクタ間メッセージ\\
\verb|.../RotLines.abcl| & ABCL + WM & アニメーション\\
\bottomrule
\end{tabular}
\end{center}

\subsection{hello.c / sum.c}
\begin{code}
int main(void) { printf("hello...\n"); return 0; }
\end{code}
\verb|#include| は受理して無視、\verb|printf| はビルトイン 0、戻り値は
\verb|aoutRun()| の戻り値になります。簡略バイトコード:
\begin{code}
ENTER 0                         ; main は局所なし
PUSH_STR offset_of("hello...\n")
CALL_BI 0, 1                    ; printf, 引数 1
PUSH_I32 0
RET
\end{code}
\verb|sum.c| は \verb|for| ループの例 (cc は \verb|for(init;cond;step)| を
\verb|init; while(cond){body; step;}| に展開。\verb|+= ++ --| は未サポート):
\begin{code}
xsh$ cc /home/sum.c -o /home/sum
xsh$ run /home/sum
sum 1..10 = 55
\end{code}

\subsection{PingPong.abcl --- アクタ間メッセージング}
\begin{code}
class Player {
    var hits;
    method bounce(other, n) {
        if (n > 0) {
            printf("Player %d: tick (n=%d) hits=%d\n", self, n, hits);
            hits = hits + 1;
            send other.bounce(self, n - 1);
        }
    }
}
main { new Player p1; new Player p2; send p1.bounce(p2, 6); }
\end{code}
\verb|self| はメソッド内の現在アクタ ID (生成順 0,1,\dots)、\verb|send| は
\textbf{非同期}、\verb|n-1| が 0 で \verb|if| を抜けると \verb|send| が出ず連鎖は
自然停止します。
\begin{code}
Player 0: tick (n=6) hits=0
Player 1: tick (n=5) hits=0
Player 0: tick (n=4) hits=1
...
Player 1: tick (n=1) hits=2
\end{code}

\subsection{RotLines.abcl --- ABCL + WM のアニメーション}
\begin{code}
class Rotator {
    var angle;
    method spin(n) {
        if (n > 0) {
            wm_line(0, 512, 384,
                    512 + icos(angle) * 320 / 4096,
                    384 + isin(angle) * 320 / 4096, rgb(255, 80, 80));
            // スロット 1..3 を +90/+180/+270 度で同様に
            sleep_ms(16);
            angle = angle + 3;
            send self.spin(n - 1);
        } else { wm_render(0); }
    }
    method start(n) { wm_render(1); send self.spin(n); }
}
main { new Rotator r; send r.start(360); }
\end{code}
\verb|icos|/\verb|isin| は Q12 固定小数を返すので半径を掛けたら 4096 で割る。
\verb|sleep_ms(16)| でだいたい 60 fps。\verb|self| への再帰送信でアニメーション
ループを作り、終了条件 (\verb|n == 0|) で \verb|wm_render(0)| を呼んで線レイヤを
消します。

\subsection{rotlines (シェルコマンド) と apps/abcl\_program.c}
\verb|shell/xsh_rotlines.c| は \verb|rot_render()| を \verb|wm_set_user_render()|
で WM に渡し、メインスレッドは角度を進めながら \verb|sleep(16)| を繰り返すだけ。
描画は (スロット型の \verb|wm_line| ではなく) 毎フレーム関数の \verb|wm_draw_line|。
\verb|apps/abcl_program.c| は「abclc が ABCL から\textbf{生成する側}の C」の見本で、
有限バッファ問題 (\verb|Buffer|/\verb|Producer|/\verb|Consumer|/\verb|Controller|
クラス、アクタ毎にメールボックス + スレッド、\verb|dispatch_*| メガディスパッチ)
を実装しています。ABCL の \verb|class| が C の \verb|dispatch_<Class>| に、
\verb|field| が \verb|enum| インデックスに、\verb|send| が \verb|enqueue| に対応
する様子を読み取れます。

%======================================================================
\section{開発ワークフロー}
\begin{code}
$ ./compile/run-window.sh
xsh$ cd /home && make && ./hello && ./sum
xsh$ edit /home/myprog.c          # C-x C-s 保存、C-x C-c 終了
xsh$ cc myprog.c -o myprog && ./myprog
xsh$ cd /home/abclcp/abclc
xsh$ edit MyActor.abcl && abclc MyActor.abcl && ./MyActor
xsh$ rotlines
\end{code}
ヒント: 1 ターゲットだけなら \verb|make MyActor|、全部作り直すなら
\verb|rm Makefile && make|。XFS の inode は 32768 までなので大量作成は注意。

%======================================================================
\section{トラブルシューティング}
\begin{center}\footnotesize
\begin{longtable}{@{}p{5.0cm} p{8.4cm}@{}}
\toprule \textbf{症状} & \textbf{原因 / 対処}\\ \midrule \endhead
\verb|qemu-system-arm not found| & \verb|brew install qemu| / \verb|apt install qemu-system-arm|\\
ビルドが古いまま起動する & \verb|compile/| で \verb|make clean && make PLATFORM=arm-qemu|\\
シェルプロンプトが出ない & \verb|Ctrl-A| \textrightarrow{} \verb|c| で QEMU モニタ、\verb|info registers|\\
\verb|cc: line N: ...| エラー & サポート外構文 (構造体・ポインタ・配列・浮動小数 等)\\
\verb|abclc: translation failed| & 1 メソッドに \verb|send| 2 個、\verb|var|/メソッド数オーバ 等\\
\verb|command not found| だが実体あり & 暗黙 \verb|run| はパス 1 トークンのみ。引数は \verb|run PATH ARGS...|\\
WM ウィンドウが出ない & \verb|arm-qemu| ビルドか? \verb|run-window.sh| を使ったか?\\
ファイルが消えた & RAMDISK は揮発性。ホスト側にバックアップを\\
パスが長すぎる & \verb|XFS_PATH_MAX = 256|。長すぎるパスは切り詰め\\
画面が崩れた & エディタ中なら \verb|C-l|、シェルなら \verb|clear|\\
\bottomrule
\end{longtable}
\end{center}

%======================================================================
\section*{付録 A. ファイルマップ}
\addcontentsline{toc}{section}{付録 A. ファイルマップ}
\begin{center}\footnotesize
\begin{tabular}{@{}l p{9.0cm}@{}}
\toprule \textbf{領域} & \textbf{主要ファイル}\\ \midrule
起動 & \verb|system/main.c|, \verb|system/initialize.c|, \verb|compile/run-*.sh|\\
XFS  & \verb|include/xfs.h|, \verb|system/xfs.c|, \verb|device/ramdisk/|\\
cc   & \verb|include/aout.h|, \verb|system/cc.c|, \verb|system/aout.c|, \verb|shell/xsh_cc.c|\\
abclc& \verb|include/abclc.h|, \verb|system/abclc.c|, \verb|shell/xsh_abclc.c|\\
make/edit & \verb|shell/xsh_make.c|, \verb|shell/xsh_edit.c|\\
WM   & \verb|apps/wm.c|, \verb|apps/abcl_xinu_gui.c|, \verb|shell/xsh_rotlines.c|\\
プラットフォーム & \verb|compile/platforms/arm-qemu/|, \verb|compile/platforms/arm-rpi/|\\
\bottomrule
\end{tabular}
\end{center}

\section*{付録 B. 既存サンプル一覧 (FS シード後)}
\addcontentsline{toc}{section}{付録 B. 既存サンプル一覧}
\begin{code}
/home/
├── hello.c
├── sum.c
└── abclcp/
    └── abclc/
        ├── PingPong.abcl
        └── RotLines.abcl
\end{code}
\verb|cd /home && make| で \verb|hello|, \verb|sum| が、
\verb|cd /home/abclcp/abclc && make| で \verb|PingPong|, \verb|RotLines| が
それぞれ a.out として生成されます。さらに深く知るには \verb|system/cc.c| と
\verb|system/aout.c| (バイトコード VM)、\verb|system/abclc.c| (ABCL 翻訳) を
読むのが近道です。

\end{document}
